\section{The projection engine, and why it is worth validating first}
\label{sec:engine}

A mechanism is only worth reporting if a reader can tell the difference between
its effect and an arithmetic error. That requirement drove the order of the work:
the accounting layer was built and validated against published outputs before any
mechanism was added to it, and the mechanism was then implemented as an extra
axis on the same code rather than as a new model.

\subsection{The accounting layer}

The engine is a standard cohort-component projection with no theory in it. It is
single year of age from 0 to 100 with an open-ended final group, two sexes, 237
countries and territories, annual steps, populations dated 1 January and rates in
force during the calendar year. It takes fertility, survival and net migration as
given and does the bookkeeping. Its inputs come from World Population Prospects
2024 \citep{un2024methodology}: age-specific fertility, survival ratios into each
age, and sex ratio at birth.

Two conventions are worth stating because getting them wrong produces plausible
rather than obviously broken output. Survival $s[a]$ is survival \emph{into} age
$a$, which is the UN's own convention, so that the numbers in the engine are the
numbers in the source file with no translation step. And childbearing ages run
from 10 to 54 rather than 15 to 49: the UN's single-age fertility file stops at
49, its five-year file does not, and the missing mothers are about 0.3\% of world
births --- small in one year, and a visible bias by 2100.

\subsection{The validation}

The test that matters is the zero-migration variant, because every input to it is
published and nothing in it is tuned, so any discrepancy is our arithmetic being
wrong. Running from 1 January \PhaseFiveBaseYear{} to 1 January 2100 on the UN's
own inputs (\Cref{tab:validation}), world population lands within
\ValidationWorld{} of the UN's published figure, the worst of
\ValidationGatedCountries{} countries above ten thousand people is within
\ValidationCountry{}, and the worst five-year age group below 100 is within
\ValidationAge{}.

\begin{table}[t]
\centering
\small
\begin{tabular}{lrr}
\toprule
Check at 2100, WPP 2024 zero-migration variant & Result & Tolerance \\
\midrule
World population                                  & \ValidationWorld{}   & 0.05\% \\
Worst country above 10{,}000 people (\ValidationWorstCountry{}) & \ValidationCountry{} & 0.5\% \\
Worst five-year age group below 100               & \ValidationAge{}     & 0.05\% \\
Open-ended 100+ group                             & \ValidationOpenAge{} & 2\% \\
Constant fertility vs the UN's own variant        & \ValidationConstantFertility{} & 0.5\% \\
\bottomrule
\end{tabular}
\caption{The engine reproduces the United Nations' published projection from the
United Nations' published inputs. The zero-migration variant is the real test
because nothing in it is estimated by us. The medium variant agrees to
\ValidationMediumDiagnostic{} but is reported as a diagnostic rather than a test,
because the UN does not publish net migrants by single year of age and ours are
backed out of its own medium path as a residual.}
\label{tab:validation}
\end{table}

The constant-fertility comparison is the absurdity check at the other end of the
range: freezing the base year's rates gives about
\ConstantFertilityTwentyOneFifty{} billion people by \ProjectionEndYear{}, within
\ValidationConstantFertility{} of the UN's own constant-fertility variant at
2100. Worth recording is that an earlier specification of this project expected
roughly 244 billion, a figure taken from a UN long-range report built on the 2002
revision. Rates have fallen a great deal since 2002, so reaching 244 billion from
a 2024 base would have meant the engine was wrong. A remembered number was
treated as a requirement for longer than it should have been.

\subsection{Why the mechanism's effect is attributable}

The typed engine of \Cref{sec:mechanism} is the same code with a composition
axis. Setting every propensity to one leaves the composition free to move but
unable to change fertility, and in that configuration the typed engine reproduces
the ordinary engine to a relative difference of $3\times10^{-16}$ --- machine
precision. It also lands on the same \ProjectionEndYear{} world total as the
older single-axis long-run diagnostic.

This is not a formality. It means the difference between a run with selection and
a run without it cannot be a second implementation of the arithmetic, a different
handling of the open age group, or a changed exposure convention. It is the
mechanism. Separately, one generation of the typed engine matches the analytic
parent--offspring covariance response of \Cref{eq:response} to nine decimal
places, which is the implementation agreeing with its own calibration.

\subsection{Migration, and what it is allowed to support}

The UN does not publish net migrants by single year of age. For the runs in
\Cref{sec:results} the migration input is backed out of the UN's own medium path
as a residual, which makes it a usable forward-model input and no evidence at
all: it absorbs every difference between the UN's procedure and ours. Runs using
it are labelled diagnostics.

Where migration uncertainty is the object of study --- \Cref{sec:horizon} and the
robustness check in \Cref{sec:boundary} --- a different and stronger construction
is used: the published Bayesian migration trajectories of \citet{azose2015}, with
rates applied to each path's own evolving population and every draw-year balanced
so that world net migration is exactly zero. That constraint matters more than it
sounds. The public trajectory archive is balanced in expectation but not draw by
draw, and using it unbalanced allows individual paths to create or delete millions
of people globally. An earlier version of the uncertainty decomposition reported
below did exactly that and attributed 1.75 billion of world spread to migration; a
component that cancels globally by construction cannot do that, and the corrected
figure is \WidthMigration{} billion.
